\chapter{RobustCluster-Aβ: robustness extension}
\section{Temporal correlation and effective sampling}
For each replica and baseline descriptor, the extension subtracts the mean and computes autocorrelation using an FFT-based convolution. Covariances use the common $1/N$ normalization. It pairs adjacent terms at positive lags and stops at the first pair whose sum is nonpositive. Each retained pair sum is then limited to the smallest pair sum seen so far. This gives the diagnostic
\begin{equation}
g=\max\!\left(1,1+2\sum_{k\in\mathcal K}\widehat\rho(k)\right),\qquad
N_{\mathrm{eff,corr}}=N/g,\qquad \tau_{\mathrm{int}}=g\Delta t/2.
\label{eq:correlation}
\end{equation}
In \cref{eq:correlation}, $\mathcal K$ denotes the lag contributions retained by this procedure; the implementation saves the lag at which summation stops. The convention includes the zero-lag half term in $\tau_{\mathrm{int}}$. Some definitions omit this term, so a reported ``correlation time'' needs an explicit convention.

Three checks help assess the estimate: whether a stopping lag was found, whether $N\geq50g$, and whether the absolute difference between first- and second-half means is at most 0.25 of the full-series standard deviation. The result is flagged if any check fails. These practical checks identify estimates with weak support from the available series; they are not a formal test that the process is stationary. More generally, an uncertainty estimate must be supported by the data and its sampling assumptions \autocite{grossfield2019,chodera2016}.

Weight concentration is assessed separately by
\begin{equation}
N_{\mathrm{eff,Kish}}=\frac{(\sum_i w_i)^2}{\sum_i w_i^2}.
\label{eq:kish}
\end{equation}
Equation~\ref{eq:kish} describes how unequal weights reduce the effective number of contributing observations, without correcting for dependence between them. It is not multiplied by, or automatically divided by, the correlation ESS. Combining the two effects would require a dependence calculation for the particular weighted observable and estimator. No Kish ESS is reported for the real data because their statistical weights have not been validated.

\section{Cluster stability across seeds}
Random grouping affects which local clusters are found, so the planned DPC seed experiment repeats discovery, consensus cutting and final assignment together. Simply renaming the labels of a fixed partition would not test this variation. The smoke profile uses two seeds, the separate pilot uses five, and the full profile, which has not been run, defaults to 20, with support for a 20--50-seed study.

Partitions are compared using the same source frames. ARI measures agreement regardless of the names assigned to clusters \autocite{hubert1985}. Families are then matched across runs to maximize their overlap, allowing Jaccard overlap, population variation and centroid movement to be summarized. Families without a match remain in the comparison, including those that appear or disappear. Considering only matched families could hide changes in cluster number even when those families look stable.

To compare centroids from different representations, the analysis uses the original four descriptors rescaled by the common baseline ranges. It therefore does not compare a PCA-space centroid directly with a raw-space centroid. Population variation across seeds measures sensitivity to the algorithm's random choices on the same ensemble. It differs from sampling uncertainty between independent replicas.

\section{Feature redundancy and PCA}
The completed lightweight diagnostics calculate the pooled Pearson correlation matrix and PCA after standardizing each baseline feature to zero mean and unit variance. The explained-variance spectrum shows how much of the variation in the observed descriptors lies along each principal component. It does not measure thermodynamic relevance or identify a slow physical mode.

The alternative clustering configurations retain all principal components, with optional whitening to scale each component to unit variance. Keeping all components makes PCA an orthogonal rotation of the standardized variables, preserving their Euclidean distances. Any difference from the min--max baseline would therefore come from standardization, whitening, numerical effects or a later analysis choice, not from discarding dimensions. This distinction matters when interpreting a future change in clustering.

\section{Independent GMM feature-ablation experiment}
The completed original experiment compares the four baseline descriptors with $\na$, $\nb$ and $R_g$ after removing SASA. It uses Gaussian mixtures in scalar descriptor space, a different model from both the published DPC pipeline and coordinate-based ShapeGMM. For a standardized observation $\mathbf{x}$,
\begin{equation}
p(\mathbf{x})=\sum_{j=1}^{K}\pi_j\mathcal{N}(\mathbf{x}\mid\symbf{\mu}_j,\symbf{\Sigma}_j),\qquad \sum_j\pi_j=1.
\label{eq:gmm}
\end{equation}
Equation~\ref{eq:gmm} approximates the sampled descriptor distribution. Its mixture coefficients are not equilibrium statistical weights. The implementation uses scikit-learn \autocite{pedregosa2011}; each component has a full covariance matrix, allowing the model to represent correlations between its features.

The experiment uses one fixed, replica-balanced fitting sample: 1,000 frames from each replica, selected with sampling seed 20261007. Feature means and standard deviations are fitted on those 3,000 rows and reused for the shared features in both representations. At each $K=2,\ldots,6$, five seeds (42--46) initialize separate models, each with a 500-iteration limit, tolerance $10^{-3}$ and covariance regularization $10^{-6}$. Using one initialization per fit exposes variation between local solutions instead of selecting the best of several starts. The protocol was committed before execution.

Every fitted model assigns a label to all 300,003 observations, including those used for fitting. ARI compares the four- and three-feature assignments at each seed and component count; the ten seed pairs are also compared within each representation. Matching clusters by maximum overlap shows which labels correspond across runs. The analysis records Jaccard overlap, unmatched families, and centroid shifts measured in the shared standardized three-feature space. It also records raw label fractions for each replica. These comparisons describe the same data used in the analysis; they do not test prediction on held-out data, and the ranges across runs are not confidence intervals.

Keeping the fitting subset fixed isolates the effects of feature representation and initialization. The experiment does not measure uncertainty from selecting that subset or remove temporal dependence. BIC is saved only for comparisons within a feature space; it is not compared across dimensions or used to identify a physical number of states. The planned DPC ablations remain separate experiments that require their own input and execution records.

\section{Wasserstein scaling sensitivity}
The raw consensus distance and the normalized alternative both combine separate distances for each descriptor; only the scaling changes. This tests the effect of one unresolved implementation choice. Raw SASA and radius-of-gyration values have different units and numerical ranges, so the variable with the larger range can have a disproportionate influence on the combined distance.

The normalized comparison would measure how scaling affects the result, but would not establish which implementation the original authors intended. Joint distances, sliced Wasserstein, energy distance and MMD could test other aspects of similarity between distributions in later studies. None replaces the current baseline.

\section{Reweighting sensitivity}
The optional comparison keeps the fitted partition fixed and compares the final-bias estimator with weights derived from verified \code{metad.rbias}. It examines the reweighted distributions of all four baseline descriptors and family relevance. With the labels unchanged, any differences can be traced to weighting rather than to changes in the partition.

The descriptor file contains a column named \code{metad.rbias}, but the original PLUMED configuration and evidence for its energy unit are missing. Input validation therefore blocks this optional comparison. The implementation keeps the estimators separate, as described in the primary literature and official PLUMED reference \autocite{tiwary2015,bonomi2009,plumedreweight}.

\section{Alternative density-based clustering}
DPA would provide a useful benchmark because it estimates density and identifies significant peaks differently \autocite{derrico2021}. The comparison should examine family number, location, population and correspondence using the same frames and validated weights. No DPA fit has been performed in this project.

Weighted ShapeGMM can analyze biased structural ensembles \autocite{sasmal2024}. Applying it here would require atomic structures, consistent atom selections and alignment, and a verified link between structures and descriptor rows. The structures cannot be recovered from scalar descriptors alone, so this comparison remains deferred.

\section{Optional external validation}
An Aβ40 ensemble constrained by experimental data could provide an independent comparison of global compaction and secondary-structure behavior \autocite{borthakur2025}. Before comparing values, the sequence, temperature, weights, atom selections and definitions of each observable must be checked. A mass-weighted $R_g$ and a Cα-only $R_g$, for example, need not have directly comparable distributions.

No external ensemble has been processed in this project. The assessment identifies the relevant study and the conditions needed for a valid comparison. Testing agreement with NMR/SAXS-constrained descriptors remains future work.
